\scanpage{170}
\mspage{119}
% archive sheet number: 170
induits par les inclusions $X_i\unsure{-|\to}X$.

La loi fonctorielle des $p_{i\,!}(e_{X_i})$ en $i$ a été explicitée en cours de démonstration. Notons que lorsque $X$ sur $I$ \deleted{\unsure{provient d'une}}\inserted{est cofibrée} et scindée, i.e. provient d'un foncteur

\[
i\longmapsto X_i,\quad I\deleted{\unsure{\to|\mathrm{Cat}}}\longrightarrow\mathrm{Cat},
\]
alors la loi fonctorielle de $p_{i\,!}(e_{X_i})$ en $i$ s'explicite de façon particulièrement \unsure{agréable} : à une flèche

\[
\alpha:i\to j,\quad\text{d'où }\alpha_*:X_i-X_j,
\]
est \unsure{associé|associée}

\[
p_!(\alpha):p_{i\,!}(e_{X_i})\longrightarrow p_{j\,!}(e_{X_j}),
\]
en notant qu'on a un diagramme

\diagram{\begin{tikzcd}[ampersand replacement=\&,column sep=small]
X_i \arrow[rr,"\alpha_*"] \arrow[dr,"p_i"',""{name=src,pos=.63}] \& \& X_j \arrow[dl,"p_j",""{name=tgt,pos=.63}]\\
\& Y \arrow[from=src,to=tgt,"\lambda_\alpha",bend left=7] \& {,}
\end{tikzcd}}
avec une flèche de commutation

\[
\lambda_\alpha:p_i\longrightarrow p_j\alpha_*
\]
(pas iso en \uline{général}), d'où

\[
p_{i\,!}(e_{X_i})\xrightarrow{(\lambda_\alpha)_!}
p_{j\,!}\underbrace{\alpha_{*\,!}(e_{X_i})}_{\longrightarrow e_{X_j}}
\longrightarrow p_{j\,!}(e_{X_j}),
\]
\unsure{plus} de doute c'est lui ! \margin{\tikz[baseline=-.2ex]{\draw[->] (0,0.65) .. controls (0,.22) and (.06,.13) .. (.52,0);} cette flèche devrait être explicitée avec soin !}
